% Section: P8 breakeven-tier analysis on iter-160 VALUE-max outputs (iter 164 JOB A)
% Fresh vein. Builds on iter-160's τ*-optimization by answering the operational
% follow-up: "at each (rate × fset), what is the CHEAPEST tier at which
% VALUE-max util recovers a target (0.99, 0.95, 0.90)?"
%
% Verdicts (6/6 PASS):
%  H1: cheap_heuristic clears target on ≥80% cells — 60/60 = 100.0%
%  H2: small_open clears target on ≥80% cells — 60/60 = 100.0%
%  H3: frontier_gpt4 recovers util ≥ 0.95 on ≥90% cells — 60/60 = 100.0%
%  H4: τ*_VALUE-max monotone in tier (cheap → frontier) on ≥80% cells — 20/20 = 100.0%
%  H4a: breakeven tier = cheap_heuristic on ≥80% cells — 60/60 = 100.0%
%  H5: frontier_gpt4 matches cheap_heuristic's target on ≥80% cells — 59/60 = 98.3%

\subsubsection{Breakeven-tier analysis on VALUE-max outputs}
\label{sec:p8-iter164-breakeven}

The iter-160 finding that VALUE-max EXCEEDS F1-max at every
(\ref{sec:p8-iter160}) cell invites the operational follow-up: \emph{at
each (rate $\times$ fset), what is the cheapest cost tier at which
VALUE-max utility recovers a chosen target level (0.99, 0.95, 0.90)?}
\textbf{Iter-164} answers this question with no additional training — the
2{,}000 iter-160 cells are re-aggregated.

\paragraph{Setup.}
For each (rate $\times$ fset $\times$ target $\in \{0.99, 0.95, 0.90\}$) cell,
we compute the per-tier mean VALUE-max util across the 5 seeds and find
the cheapest tier at which the mean clears the target. There are
$5 \times 4 \times 3 = 60$ cells. The five-tier price ladder matches
iter-148: cheap\_heuristic (\$0.0001/call), small\_open (\$0.0006/call),
iter120\_default (\$0.0010/call), mid\_tier (\$0.0050/call),
frontier\_gpt4 (\$0.0300/call).

\paragraph{Headline findings (6/6 PASS).}
\begin{enumerate}
\item \textbf{cheap\_heuristic clears every target on every cell.}
H1 60/60 = 100.0\% of (rate $\times$ fset $\times$ target) cells reach
the target at the cheapest tier. The breakeven tier is therefore
\emph{always} cheap\_heuristic — H4a 60/60 = 100.0\%.
\item \textbf{small\_open also clears every target on every cell.}
H2 60/60 = 100.0\%. The 6$\times$ per-call price premium over
cheap\_heuristic buys nothing on VALUE-max.
\item \textbf{frontier\_gpt4 still recovers $\geq 0.95$ on every cell.}
H3 60/60 = 100.0\% (min observed cell mean = 0.9885). Even at \$0.03/call
with rare-event alert volume variance, VALUE-max is dominated by
$V_{\text{per\_catch}} = \$50$.
\item \textbf{$\tau^*_{\text{VALUE-max}}$ is monotone increasing in tier.}
H4 20/20 = 100.0\% of (rate $\times$ fset) cells show
$\tau^*_{\text{cheap}} \leq \tau^*_{\text{small\_open}} \leq \cdots \leq
\tau^*_{\text{frontier}}$. The costlier tier pushes the threshold up
to suppress false alerts — at cheap\_heuristic $\tau^* \approx 0.55$
(recall-pushing); at frontier\_gpt4 $\tau^* \approx 0.99$
(precision-pushing).
\item \textbf{frontier\_gpt4 matches cheap\_heuristic on the same target.}
H5 59/60 = 98.3\% of (rate $\times$ fset $\times$ target) cells show
frontier\_gpt4 also clearing the target. The single failure cell is
0.05\% $\times$ 20raw+stat $\times$ 0.99 target, where the
frontier tier drops to 0.9885 (1.15 percentage-point degradation).
\end{enumerate}

\paragraph{Sharpest negative finding.}
The 1 of 60 cells where frontier\_gpt4 fails the 0.99 target is the
0.05\% rate (rarest-event regime) on the 20raw+stat feature set
(sparsest non-aggregated features) at the 0.99 target. At this cell,
seed 20260714 saw 248 alerts, \$7.44 cost, and \$250 value\_gain, yielding
$\mathrm{util} = 0.97024$. The other 4 seeds on this cell range from
0.9957 (36 alerts) to 0.9893 (89 alerts). Alert-volume variance
at the 0.05\% rate is the binding constraint — not the tier price.

\paragraph{Operational recommendation.}
\begin{itemize}
\item \textbf{Default tier: cheap\_heuristic (\$0.0001/call).} Clears
every reasonable utility bar on every (rate $\times$ fset) cell. The
savings relative to small\_open are 6$\times$ per call; relative to
frontier\_gpt4 are 300$\times$ per call.
\item \textbf{Do not pay for frontier LLMs to flag fraud.} VALUE-max
at frontier\_gpt4 is dominated by cheap\_heuristic on every cell —
frontier only buys redundancy, not value.
\item \textbf{$\tau^*$ at cheap\_heuristic $\approx 0.55$} (recall-pushing
VALUE-max); at frontier\_gpt4 $\approx 0.99$ (precision-pushing).
\end{itemize}

\paragraph{Cross-paper coupling.}
This finding sharpens iter-160's VALUE-max dominance: iter-160 showed
VALUE-max EXCEEDS F1-max at every cell; iter-164 shows VALUE-max clears
operational bars (0.99, 0.95, 0.90) at the cheapest tier on every cell.
The two findings together mean \emph{VALUE-max is the dominant utility
function at every tier} — the cost-tier choice is decoupled from the
utility-function choice.